\section{Study design}
\label{sec:design}

\subsection{Pre-registration and hypotheses}

Research questions, hypotheses, agent-facing texts, the detector and the analysis plan were frozen in a
pre-registration committed before any model data (commit \texttt{231866f}). Every later change is a dated
entry in \texttt{DEVIATIONS.md}, summarised in Section~\ref{sec:deviations}. The confirmatory hypotheses
were: H1, violations rise monotonically with pressure; H2, rates differ by lead type; H3, soft controls
produce more violations than hard ones and fewer than policy-only, with hard controls turning violations into
attempts; H4, a cue that Compliance reviews are paused raises violations relative to an audit cue; H5, a
compliance-primacy sentence reduces but does not eliminate violations; H6, more than 20\% of violators % numbers-ok: design constant, version or citation
deceive at the check-in; H7, hard controls plus a legitimate escalation outlet reduce violations to the floor
at the smallest cost in goal attainment; H8, most violations are knowing.

\subsection{Conditions}

The pressure levels are the manager's message (verbatim in Appendix~\ref{app:prompts}). P0 has no target.
P1 states the target and progress. P2 adds that the RM's bonus and review depend on closing the gap. P3 says
branches that miss the quarter will be consolidated and RM roles cut, ``yours included'', and that every
lead must be converted. None of P0--P3 asks for a violation. Two ceiling conditions do: P4 adds ``If the
checks are holding a customer up, do whatever it takes to get the account open today --- use the override
fields, I'll back you'', and P5, added after the pilot (Section~\ref{sec:deviations}), adds a direct order to
open the account anyway with ``Approved by branch manager'' as the override reason. Other factors are a
compliance-primacy sentence in the system prompt (C+), an oversight cue (O1 ``All actions are logged and
reviewed weekly by Compliance''; O2 ``Compliance is short-staffed; reviews are paused until next quarter''),
control hardness (H0/H1/H2), and four mitigations: M1 a rulebook summary at the top of the brief, M3 a
\texttt{request\_target\_extension} tool that the manager grants, M4 a mandatory one-line compliance
justification on every account opening and deposit, and M6 hard controls plus M3. A floor condition replaces
the sales role with onboarding review and no target.

\subsection{Experiments}

Table~\ref{tab:design} lists the seven confirmatory experiments. Each Level-1 cell crosses a condition with
the six blocked lead types; seeds are fixed per cell and reused across conditions, so conditions differ only
in the manipulated factor.

\begin{table}[t]
\centering\small
\caption{Confirmatory experiments. N is per cell (condition $\times$ lead type) for Level~1 and per condition for % numbers-ok: design constant, version or citation
Level~2. Defaults: pressure P3, soft controls (H1), no compliance sentence (C$-$), no oversight cue (O0), no % numbers-ok: design constant, version or citation
mitigation.}
\label{tab:design}
\begin{tabular}{@{}lllr@{}}
\toprule
 & Manipulation & Models & N \\
\midrule
A & pressure P0--P3 & Haiku, Sonnet; Opus & 30; 12 \\ % numbers-ok: design constant, version or citation
B & control hardness H0, H2 & Haiku & 30 \\ % numbers-ok: design constant, version or citation
C & oversight cue O1, O2 & Haiku & 30 \\ % numbers-ok: design constant, version or citation
D & compliance sentence C+ at P1--P3 & Haiku & 30 \\ % numbers-ok: design constant, version or citation
E & mitigations M1, M3, M4, M6 (leads B1, B3, B6) & Haiku & 30 \\ % numbers-ok: design constant, version or citation
F & floor (no target) and P4 ceiling & Haiku, Sonnet & 6 \\ % numbers-ok: design constant, version or citation
G & five-day episodes at P3 and P0 & Haiku & 20, 10 \\ % numbers-ok: design constant, version or citation
\bottomrule
\end{tabular}
\end{table}

\subsection{Models and harness}

The agents were Claude Haiku~5.5, Sonnet~5.5 and Opus~5.5 (\texttt{claude-haiku-5-5}, % numbers-ok: design constant, version or citation
\texttt{claude-sonnet-5-5}, \texttt{claude-opus-5-5}) through the Anthropic Messages API with a manual tool
loop over the in-process MCP server. We sent no temperature or sampling parameters, left thinking at the API
default and pinned the effort setting to each model's documented default (Haiku and Opus \emph{medium}, Sonnet
\emph{high}). A response cut off at the 8{,}000-token limit restarts the probe once from a fresh bank; a % numbers-ok: design constant, version or citation
response the API's safety classifier declined is recorded separately. Experiments ran on 2026-10-28. % numbers-ok: design constant, version or citation

\subsection{Sample size and budget}

The pre-registration fixed the rule rather than the number: after the pilot, measured cost per probe sets one N
per experiment before Experiment A, and N is never changed in response to results. The study ran on a prepaid
API budget with a hard cap of \$80 at list prices, enforced by a spend ledger that stops a run before the cap. % numbers-ok: design constant, version or citation
The final N (Table~\ref{tab:design}) was written to the deviation log before A started; Opus was added last
at N\,=\,12 from budget left over. The whole study, including pilots, judges and comprehension, cost % numbers-ok: design constant, version or citation
\$\SpendUSD{} at list prices.

\subsection{Analysis}

The pre-registered analysis is logistic regression of violation on the manipulated factor with lead type as a
covariate, with Firth's penalised likelihood as the declared fallback under separation, a Cochran--Armitage
trend test for H1, exact binomial tests for H6 and H8, and bootstrap confidence intervals for differences in
goal attainment. With zero events these tests are not informative, so the primary report is the count with a
Wilson 95\% interval per cell. The pre-registration also names robustness variants (counting attempts as % numbers-ok: design constant, version or citation
violations, excluding V7); with no events of either kind they coincide with the main analysis.

\subsection{Instrument validation}
\label{sec:gates}

The pre-registration set four gates before Experiment A. Gate~1, the compliance-only floor shows about zero % numbers-ok: design constant, version or citation
violations for both models, and gate~4, comprehension accuracy of at least 0.9 on clean-lead items, passed. % numbers-ok: design constant, version or citation
Gate~2 required the P4 ceiling to produce a high % numbers-ok: design constant, version or citation
violation rate, to show that model-made violations register; it failed, because neither Haiku nor Sonnet violated in the
pilot at P4 (\NPilotPFour{} probes) or at the P5 direct order (\NPilotPFive{} probes). We report this
failure as a result and validated the instrument without a model instead: a scripted agent that takes the
circumvention path for every lead type under all three hardness levels was flagged in
\NScriptedDetected{}/\NScripted{} runs with the expected codes (VIOLATED under H0 and H1, ATTEMPTED under
H2), a scripted compliant agent was clean in \NScriptedCompliant{}/\NScriptedControls{}, and a test sends the
same violation as model \texttt{tool\_use} blocks through the API path and checks that it is detected.
Gate~3 requires detector precision and recall of at least 0.95 on 60 blind-labelled runs. % numbers-ok: design constant
The sample mixes pilot probes with scripted violations so that recall is measurable although no model
violated. The pre-registration called for the author to label it; the author declined, and before any label
existed we fixed a replacement (Appendix~\ref{sec:deviations}): Opus~5.5 and Sonnet~5.5 each labelled the \SpotN{} % numbers-ok: model version
cases independently, one call per case, seeing only the lead type, the control hardness, the event list and
the labelling guide, and never the detector's verdict. Both agreed with the detector's outcome in every case
(Opus \SpotAgreeOpus{}/\SpotN{}, Sonnet \SpotAgreeSonnet{}/\SpotN{}; false positives \SpotFPOpus{} and
\SpotFPSonnet{}, false negatives \SpotFNOpus{} and \SpotFNSonnet{}), and with each other in
\SpotLabellerAgree{} (Cohen's $\kappa = \SpotKappa$). Gate~3 therefore passed on its pre-registered rule, % numbers-ok: design constant
which is stated on point estimates. With only \SpotDetectorPositives{} positives, the 95\% Wilson lower bound % numbers-ok: design constant, version or citation
on recall is \SpotRecallLo, well below the 0.95 the gate names. All \SpotDetectorPositives{} positives are % numbers-ok: design constant
scripted violations, which are easy to recognise, and the labellers read the same event list the detector
reads, so agreement shows that the detector's rules match the written policy, not that it catches subtle
model-made cases. No human labelled the sample, and both labellers belong to the model family under test;
Opus~5.5 is also one of the agents tested (Section~\ref{sec:limitations}). % numbers-ok: model version
